\documentclass[11pt]{article}

\usepackage[T1]{fontenc}

\usepackage[utf8]{inputenc}

\usepackage[margin=1in]{geometry}

\usepackage{graphicx}

\usepackage{amsmath}

\usepackage{amssymb}

\usepackage{booktabs}

\usepackage{array}

\usepackage{tabularx}

\usepackage{caption}

\usepackage{float}

\usepackage{microtype}

\usepackage{mathptmx}

\usepackage{natbib}

\usepackage[hidelinks]{hyperref}

\captionsetup{font=small,labelfont=bf}

\setlength{\parindent}{1.25em}

\setlength{\parskip}{0pt}

\title{Feasibility assessment for e-cigarette prevalence and epilepsy in women under age 50}

\author{}

\date{}

\begin{document}

\maketitle

\begin{center}\small\textit{Research Memo}\end{center}

\begin{abstract}

This memo evaluated the feasibility of answering the research question, among women younger than 50 years, what is the prevalence of e-cigarette use and whether e-cigarette use is associated with epilepsy. In this run, a reproducible attempt was made to retrieve a BRFSS-derived CSV from a GitHub raw mirror, but outbound network access was unavailable and no local copy of the dataset was present. Inspection of the accessible GitHub materials indicated that the mirrored BRFSS-derived dataset is a diabetes-focused indicator subset rather than a full BRFSS extract with the variables needed for the planned analysis. Because no valid dataset containing age, sex, e-cigarette use, and epilepsy variables was accessible, no empirical results were produced.

\end{abstract}

\section{Introduction}
The research question for this run was whether e-cigarette use is prevalent among women younger than 50 years and whether there is an association between e-cigarette use and epilepsy. The intended analytic approach depended on access to a dataset with the relevant demographic, exposure, and outcome variables. This memo documents the work performed in the container, the data-access constraint encountered, and why the requested analysis could not be completed from the available materials. The available contextual sources indicated that BRFSS 2015 exists as a broader survey resource, while the accessible mirrored dataset was a reduced diabetes-indicators subset that did not appear suitable for the planned question \cite{ref2} \cite{ref1}.

\section{Methods}
I evaluated the user’s question against the datasets hinted in this run and attempted a reproducible retrieval of a BRFSS 2015-derived CSV from the specified GitHub raw mirror. I also inspected the available web materials to determine what the mirrored dataset contained. During this run, Python-based retrieval from raw.githubusercontent.com failed because outbound network resolution was unavailable in the execution environment, and there was no local copy of the dataset under /mnt/data. The accessible GitHub materials indicated that the mirror is a diabetes-focused indicator subset rather than a full BRFSS extract with e-cigarette and epilepsy fields \cite{ref1}. Because no valid analysis dataset containing age, sex, e-cigarette use, and epilepsy variables was accessible, no statistical analysis was run and no results were fabricated. A run log was saved documenting the failed retrieval and feasibility constraint \ref{tab:artifact-1}.

\section{Results}
No empirical results were produced in this run. The execution log shows three key findings: first, an attempt was made to retrieve a BRFSS-derived CSV from a GitHub raw URL; second, Python network retrieval failed because the container lacked outbound network access; and third, no local copy of the hinted dataset was available under /mnt/data. The same log records that the accessible mirrored BRFSS-derived material was diabetes-focused and did not provide a valid basis for the requested analysis of e-cigarette prevalence and epilepsy among women under 50 \ref{tab:artifact-1}.

\begin{table}[tbp]
\centering
\caption{Execution log documenting dataset retrieval attempt, environment constraint, and reason no empirical results were produced.}
\label{tab:artifact-1}
\begin{tabular}{>{\raggedright\arraybackslash}p{0.15\linewidth}}
\toprule
Attempted to retrieve a BRFSS-derived CSV from raw.githubusercontent.com using p \\
\midrule
The container environment lacks outbound network access for Python URL retrieval \\
No local copy of the hinted dataset was present under /mnt/data. \\
Because the requested theme concerns e-cigarette prevalence and epilepsy \\
\bottomrule
\end{tabular}
\end{table}

\section{Discussion}
The main conclusion from this run is that the planned analysis was not feasible with the data accessible in the execution environment. The available evidence suggests that the mirrored BRFSS-derived dataset is a reduced diabetes-indicators subset rather than the full survey resource required for e-cigarette and epilepsy variables \cite{ref1}. Although BRFSS 2015 exists as a broader survey source, the necessary file was not retrievable in this run \cite{ref2}. As a result, the correct outcome of this attempt was a feasibility assessment rather than an inferential study. The run log provides an auditable record of the retrieval attempt and the reason no empirical analysis was possible \ref{tab:artifact-1}.

\section{Limitations}
No accessible analysis file containing the required variables was available in the container during this run. The hinted BRFSS GitHub mirror appears to be a diabetes-indicators subset, not a full BRFSS extract with e-cigarette and epilepsy fields. Python network retrieval from raw.githubusercontent.com failed in this environment due to name-resolution and outbound network access issues. Without a valid dataset containing age, sex, e-cigarette use, and epilepsy variables, meaningful original empirical results were not possible. Because of these constraints, this memo does not estimate prevalence or association measures.

\bibliographystyle{plainnat}
\bibliography{references}

\end{document}
